\chapter{LuaSTG自带的数学变量与函数} \label{sec:lstg-math}

\section{运算符}

LuaSTG使用的lua版本为luajit,
支持的运算符与lua5.1相同.

来源: 网络

\begin{center}
  \begin{xltabular}{\textwidth}{XX}
    \toprule
    运算符           & 说明                           \\
    \midrule
    {
      $a+b$, $a-b$, $a*b$, $-a$
    }                & {
        加法, 减法, 乘法, 相反数
      }                                                 \\
    $a\mathbin{/}b$  & 浮点数除法                     \\
    $a\mathbin{\%}b$ & {
        $a$被$b$除的余数, 计算结果在0到$b$之间
      }                                                 \\
    {
    $a\mathbin{\text{\textasciicircum}}b$
    }                & {
        $a$的$b$次方, 即$a^b$
      }                                                 \\
    {
    $a>b$, $a<b$, $a>=b$, $a<=b$, $a==b$,
    $a\sim=b$
    }                & 比较两个数的大小               \\
    {
    $A\mathbin{\text{and}}B$,
    $A\mathbin\text{or}B$,
    $\mathbin{\text{not}}A$
    }                & 逻辑运算符                     \\
    $\#A$            & {
        表$A$的数组长度 (只支持连续索引的数组)
      }                                                 \\
    $A\mathbin{..}B$ & 拼接两个字符串, 返回新的字符串 \\
    \bottomrule
  \end{xltabular}
\end{center}

\section{全局变量与函数}

来源: \raw{doc/core}, \raw{lib/Lapi.lua},
\raw{lib/Lmath.lua}, \raw{lib/Lscreen.lua}

\begin{center}
  \begin{xltabular}{\textwidth}{XX}
    \toprule
    变量/函数                                  &
    说明                                         \\

    \midrule
    PI                                         &
    圆周率$\pi$                                  \\

    SQRT2, SQRT3                               &
    $\sqrt2, \sqrt3$                             \\

    GOLD                                       &
    $360\degree$与黄金比例的乘积,
    约为$137.5\degree$                           \\

    int($x$)                                   &
    向下取整                                     \\

    abs($x$)                                 &
    绝对值$|x|$                                  \\

    sign($x$) &
    计算$x$的正负性, 返回$1,-1$或者$0$ \\

    max($a$, $b$, \dots), min($a$, $b$, \dots) &
    最大值和最小值                               \\

    sqrt($x$)                                  &
    平方根$\sqrt x$                              \\

    sin($x$), cos($x$), tan($x$)               &
    三角函数                                     \\

    asin($x$), acos($x$), atan($x$)            &
    反三角函数$\arcsin,\arccos,\arctan$          \\

    atan2($y$, $x$)                            &
    计算由原点指向$(x,y)$的方位角,
    $-180 < \cdot \le 180$           \\

    hypot($x$, $y$) &
    点$(x,y)$到原点的距离 \\

    mod($a$, $b$)                              &
    C风格的取余运算, $a$除以$b$的商向零取整
    (截断取整) 所得的余数                        \\

    BoxCheck(unit, left, right, bottom, top)   &
    检测unit对象是否在
    left,right,bottom,top
    四个数确定的矩形区域内                       \\

    {
    Angle($x_1$, $y_1$, $x_2$, $y_2$) \par
    或 Angle(unit1, unit2) \par
    或 Angle(unit1, $x_2$, $y_2$) \par
    或 Angle($x_1$, $y_1$, unit2)
    }                                          &
    计算由$(x_1,y_1)$指向$(x_2,y_2)$的方位角,
    角度范围与atan2相同                          \\

    Dist($x_1$, $y_1$, $x_2$, $y_2$)           &
    计算点$(x_1,y_1)$与$(x_2,y_2)$的距离.
    与Angle一样还有其他参数格式                  \\

    WorldToUI($x$, $y$) &
    从world坐标系 (\raw{SetViewMode("world")})
    到ui坐标系 (\raw(SetViewMode("ui")))
    的坐标转换 \\

    WorldToScreen($x$, $y$) &
    从world坐标系到屏幕坐标系的坐标转换 \\
    \bottomrule
  \end{xltabular}
\end{center}

\section{随机数生成}

为了与replay系统适配, LuaSTG提供了一个随机数生成器ran,
下面是它提供的随机数生成函数.

如果某些情况需要你另外创建一个随机数生成器,
可以使用Rand()函数进行创建.

\begin{center}
  \begin{xltabular}{\textwidth}{XX}
    \toprule
    函数 & 说明 \\
    \midrule

    ran:Int($a$, $b$)   &
    生成$a$到$b$范围的整数, 包含$a$和$b$ \\

    ran:Float($a$, $b$) &
    生成$a$到$b$范围的浮点数             \\

    ran:Sign()          &
    生成$-1$或1                          \\
    \bottomrule
  \end{xltabular}
\end{center}

\section{math库}

Lua语言自带的math库提供了一些数学运算函数.

来源: Lua5.1参考手册

\begin{center}
  \begin{xltabular}{\textwidth}{XX}
    \toprule
    变量/函数 & 说明 \\
    \midrule

    math.pi &
    圆周率 \\

    math.sqrt($x$) &
    平方根 \\

    math.abs($x$) &
    绝对值 \\

    math.floor($x$) &
    向下取整 \\

    math.ceil($x$) &
    向上取整 \\

    math.max($a$, $b$, \dots), math.min($a$, $b$, \dots) &
    最大值和最小值 \\

    math.deg($x$) &
    将弧度值$x$转换为角度值 \\

    math.rad($x$) &
    将角度值$x$转换为弧度值 \\

    \st{math.random($m$, $n$)} &
    与replay系统不兼容, 建议使用ran随机数生成器提供的方法 \\

    math.sin($x$), math.cos($x$), math.tan($x$) &
    三角函数, 传入参数为弧度值 \\

    math.asin($x$), math.acos($x$), math.atan($x$) &
    反三角函数, 返回的是弧度值 \\

    math.atan2($y$, $x$) &
    计算由原点指向$(x,y)$的方位角, 返回的是弧度值 \\

    math.pow(x,y) &
    计算$x^y$, 即运算符\textasciicircum \\

    math.exp($x$) &
    计算$e^x$ ($e$为自然对数的底) \\

    math.log($x$, base) &
    以base为底的$x$的对数 \\

    math.log10($x$) &
    以10为底$x$的对数 \\

    math.sinh($x$), math.cosh($x$), math.tanh($x$) &
    双曲函数 \\

    math.modf(x) &
    返回$x$的整数部分和小数部分 \\

    math.fmod($a$, $b$) &
    C风格的取余运算, $a$除以$b$的商向零取整
    (截断取整) 所得的余数 \\

    math.frexp($x$) &
    返回两个数$m,e$使得$x = m\cdot2^e$,
    其中$e$为整数, $m$的绝对值在0.5到1之间 \\

    math.ldexp($m$, $e$) &
    返回$m\cdot2^e$ \\
    \bottomrule
  \end{xltabular}
\end{center}

\section{位运算}

Luajit提供了对32位整数进行位运算的bit库.

来源: \raw{doc/luajit extension/bit.lua}

\begin{center}
  \begin{xltabular}{\textwidth}{XX}
    \toprule
    函数 & 说明 \\
    \midrule

    bit.tohex($x$, $n$) &
    把$x$转换为长度为$|n|$的十六进制字符串,
    $n>0$时采用小写字母, 否则采用大写字母 \\

    bit.bnot($x$) &
    按位取反 \\

    bit.band($x_1$, $x_2$, \dots) &
    按位与 \\

    bit.bor($x_1$, $x_2$, \dots) &
    按位或 \\

    bit.bxor($x_1$, $x_2$, \dots) &
    按位异或 \\

    bit.lshift($x$, $n$) &
    按位左移 \\

    bit.rshift($x$, $n$) &
    逻辑右移 \\

    bit.arshift($x$, $n$) &
    算术右移 \\

    bit.rol($x$, $n$), bit.ror($x$, $n$) &
    按位左旋/右旋 \\

    bit.bswap($x$) &
    字节翻转 \\
    \bottomrule
  \end{xltabular}
\end{center}

\section{spmath库}

LuaSTG的data层提供了一个spmath库,
详见\raw{sp/spmath.lua}.
这里省略一些用处很小或者有bug的函数.

\begin{center}
  \begin{xltabular}{\textwidth}{XX}
    \toprule
    函数 & 说明 \\
    \midrule

    sp.math.subAngle($\alpha$, $\beta$) &
    计算$\beta-\alpha$角度差,
    返回值范围为$-180<\cdot\le180$ \\

    sp.math.sum($A$) &
    数组$A$各元素的和 \\

    sp.math.circle($r$, $\alpha$, $x=0$, $y=0$) &
    与$(x,y)$距离为$r$, 方位角为$\alpha$的坐标
    (返回两个数值) \\

    sp.math.vecmul($A$, $B$) &
    向量的内积, $A,B$是两个数组 \\

    sp.math.shuffle($A$) &
    矩阵的转置, $A$是一个二维数组 \\

    sp.math.matmul($A$, $B$, toCol1, toCol2) &
    矩阵乘法, $A,B$为二维数组, 计算方法为:
    若toCol1为真, 将$A$转置,
    若toCol2为真, 将$B$转置;
    对$A$中的各个数组与$B$中的各个数组计算内积,
    得到乘积矩阵的各个元素.
    也就是说, 如果$A,B$按照行主序存储,
    那么参数应该这样填写:
    sp.math.matmul($A$, $B$, false, true) \\

    sp.math.Per3D($x$, $y$, $z$, dist) &
    伪透视计算, $x,y,z$是要投影的坐标,
    dist是摄像机与坐标原点的距离,
    计算原理如下图所示. \\
    \bottomrule
  \end{xltabular}
\end{center}

\begin{center}
  \begin{tikzpicture}
    \draw[-latex] (0,0) -- (1,0) node[right] {$z$};
    \draw[-latex] (0,0) -- (0,1) node[left] {$x$};

    \draw[dashed] (0,0) -- (-4,0)
    coordinate (0) node[left] {摄像机}
    node[midway,below] {dist};
    \draw (0) -- (2.5,1) node[right] {$(x,y,z)$};

    \draw (-4+1,0) arc (0:9:1)
    node[right,background] {$\alpha_x$};

    \draw (0) -- (0,2) node[right] {$\alpha_0$};
  \end{tikzpicture}
  \begin{align*}
    \alpha_0 = \arctan(192/\text{dist}), \\
    x' = 192\cdot \alpha_x / \alpha_0, \\
    y' = 192\cdot \alpha_y / \alpha_0.
  \end{align*}
\end{center}
